\section{The target and the structure of the argument}\label{sec:overview}

Write \(V_n=\{0,\ldots,n-1\}^2\). A set \(S\subseteq V_n\) has the
\emph{no-three-in-line property} if every three distinct points of \(S\)
are noncollinear. Equivalently,
\[
 (q_x-p_x)(r_y-p_y)-(q_y-p_y)(r_x-p_x)\ne0
 \qquad(p,q,r\in S\text{ distinct}).
\]
Every such set has at most two points on each horizontal row, hence
\(|S|\le2n\). The intended numerical conclusion is
\begin{equation}\label{eq:target}
 \exists n_0\ \forall n\ge n_0\ \exists S\subseteq V_n:
 \quad \NoThree(S),\qquad 1501n\le1000|S|.
\end{equation}
The positive improvement over \(3n/2\) is the fixed number \(n/1000\).
The quantifier is over every sufficiently large integer \(n\), with no
divisibility restriction.

\subsection{Finite input, fixed before the grid size}
Choose a finite rational time partition
\(0=t_0<t_1<\cdots<t_J=T\) and a fixed finite rational rectangular
partition of the unit square. A common refinement may be taken to be a
\(G\times G\) grid for a fixed integer \(G\). Let \(\mathcal B\) be the
union of its horizontal and vertical grid lines.
There are four parity types \(r\in(\Z/2\Z)^2\). For each type choose
a deterministic function
\[
 f_{r,t}:[0,1]^2\longrightarrow[0,1]
\]
that is constant on every open spatial cell and every time interval.
Set it to zero on \(\mathcal B\), outside the square, and after time \(T\).
Define
\begin{equation}\label{eq:input}
 m_t(p)=f_{\typ(p),t}(p/n),\qquad
 F_{r,t}(x)=\int_0^t f_{r,u}(x)\dd u,\qquad
 A=\frac14\sum_r\int_{[0,1]^2}F_{r,T}(x)\dd x.
\end{equation}
Here \(\typ(p)=p\bmod2\). The factor \(1/4\) belongs to the spatial
mass, not to the birth rate of an individual point.
The last-use time \(\tau_r(x)\) is the right endpoint of the last interval
on which that type and cell is used; it is zero for an unused cell.
For a grid point write \(\tau_p=\tau_{\typ(p)}(p/n)\).

In the candidate data the spatial mask in a phase is obtained from a
union \(U\) of accepted closed rectangles and their prescribed images:
the actual region is \(\operatorname{int}(U)\setminus\mathcal B\).
Areas are union areas. Overlapping images are never added with
multiplicity. Removing \(\mathcal B\) is part of the definition of the
discrete process, even though it has zero planar area.

\subsection{The geometric input conditions}
A primitive integer direction \(v\) is identified with \(-v\). Put
\(h(v)=\max(|v_1|,|v_2|)\) and \(\bar v=v\bmod2\ne0\).
For a field \(F\) define the major-coordinate chord integral
\[
 C_vF(x)=\int_{\R}F\bigl(x+s v/h(v)\bigr)\dd s.
\]
This is not an arc-length integral. Let \(\Gamma\) be a fixed finite
family of designated directions containing both coordinate directions,
of cardinality \(g\).
The numerical example uses all directions of height at most three,
so \(g=16\). Define the line capacity
\begin{equation}\label{eq:capacity}
 L_{v,r}(F;x)=\frac{C_v(F_r+F_{r+\bar v})(x)}{2h(v)}.
\end{equation}
Let the four direction families be
\((1,u),(1,-u),(u,1),(-u,1)\), \(0\le u\le1\), and let
\(c_{F,j}(x,u)\) denote the corresponding chord. Set
\begin{equation}\label{eq:risk}
 g_F(x)=\frac14\sum_{j=1}^4\int_0^1c_{F,j}(x,u)^2\dd u,
 \qquad
 \mathcal R_r(F;x)=\frac{D_*}{6}\sum_{s\ne r}g_{F_r+F_s}(x).
\end{equation}
Section~\ref{sec:sieve} gives a finite arithmetic criterion for a valid
coefficient \(D_*\); the example uses \(76/125\).

The finite schedule is \emph{admissible} for fixed parameters if there
are positive constants \(\kappa_H,\kappa_L\) and a number \(0\le r_*<d\)
such that
\begin{align}
 L_{v,r}(F_T;x)&\le K-\kappa_H &&(v\in\Gamma),\label{eq:heavyinput}\\
 L_{v,r}(F_T;x)&\le\lambda-\kappa_L &&(v\notin\Gamma),\label{eq:lightinput}\\
 \mathcal R_r(F_{\tau_r(x)};x)&\le r_* &&(\tau_r(x)>0),\label{eq:riskinput}
\end{align}
uniformly over the relevant roots, types and lines. In the last formula
\emph{every one of the four fields is evaluated at the same root's
deadline}. The constants satisfy
\begin{equation}\label{eq:exponents}
 \begin{gathered}
 0<K<2,\quad 0<\rho<1,\quad 0\le\lambda<1-\rho,\\
 0<d<\gamma<1,\quad 0<\xi<d,\quad \zeta>0,\quad 1-d-2\xi-\zeta>0.
 \end{gathered}
\end{equation}
All data and all displayed constants are fixed before \(n\) tends to
infinity. Constants depending on a moment order may depend on that
order, but the order is also fixed before taking the limit.

\begin{theorem}[Analytic reduction, in the scope of this proof draft]\label{thm:bridge}
For an admissible finite schedule, the stopped compensated process
defined in Section~\ref{sec:process} produces a random set \(S_n\subseteq
V_n\) with the no-three-in-line property and
\[
 \E|S_n|=(A-o(1))n.
\]
Consequently, for every fixed \(c<A\), a set of size at least \(cn\)
exists for every sufficiently large \(n\).
\end{theorem}

The detailed argument for this theorem occupies
Sections~\ref{sec:geometry}--\ref{sec:closing}. The numerical use of the
theorem is a separate obligation: the particular finite schedule must
actually satisfy \eqref{eq:heavyinput}--\eqref{eq:riskinput}, with the
claimed union mass and chronological matching. This manuscript does
not certify that entire concrete input. In particular, the theorem
is not an unconditional assertion of \eqref{eq:target}.

\subsection{Why saved credit is the useful quantity}
Each active point has a weight \(w_p\), initially one. It is selected at
rate \(m_t(p)w_p/n\). Weights grow to compensate for the chance of
selection, deletion by a newly completed pair, or a prescribed
halving. At an early safety stop, the positive weight is saved rather
than discarded. The resulting process \(Z_p\) is a bounded martingale
with \(\E Z_p(t)=1\).

The change of probability measure
\[
 \dd\Prob^{(p)}=Z_p(T)\dd\Prob
\]
measures the amount of credit lost by retiring \(p\), rather than just
the raw probability that \(p\) retires. Let \(\theta_p\) be its first
positive saving time and \(E_p\) the event that this happens strictly
before \(\tau_p\) for a safety reason. Put
\[
 A_n=\frac1{n^2}\sum_{p\in V_n}F_{\typ(p),T}(p/n).
\]
The exact identity proved in Section~\ref{sec:loss} is
\begin{equation}\label{eq:previewloss}
 A_n-\frac{\E|S_n|}{n}
 =\frac1{n^2}\sum_p\E^{(p)}\left[
 \ind_{E_p}\int_{\theta_p}^{T}m_t(p)\dd t\right].
\end{equation}
Thus a bound \(\sup_p\Prob^{(p)}(E_p)\le\varepsilon\) implies
\(\E|S_n|/n\ge(1-\varepsilon)A_n\).
There are five causes to control: weight cap, ordinary-line tag,
pair-correction floor, virtual-deletion bound, and designated-line clock.
The last three are global safety stops. The paper proves their estimates
on the already stopped process; it does not assume that those stops
are rare while deriving the estimates that make them rare.

\subsection{A fixed budget suffices for the numerical target}
The recorded scaled mass and target are
\begin{equation}\label{eq:massnumbers}
 A_* =\frac{5036868114751900569}{3355443200000000000},\qquad
 c=\frac{1501}{1000},\qquad
 \Delta=A_*-c=\frac{347871551900569}{3355443200000000000}.
\end{equation}
These are candidate input values; verifying the mass of the accepted
union remains part of the concrete application. If they are justified,
and if \(A_n\ge A_*-\Delta/4\) and
\(\sup_p\Prob^{(p)}(E_p)\le10^{-5}\), then
\[
 (1-10^{-5})(A_*-\Delta/4)-c
 =0.0000627445734908543514408171176910400390625>0.
\]
This strict rational inequality is enough for \eqref{eq:target}.
There is no need to prove an effective, practically small threshold
\(n_0\), or to compute an unbounded sequence of ever tighter error
budgets. The asymptotic argument is a way to obtain one fixed budget.

